\documentclass[11pt,a4paper]{article}
\usepackage[margin=1in]{geometry}
\usepackage{listings}
\usepackage{xcolor}
\usepackage{hyperref}
\usepackage{enumitem}
\usepackage{booktabs}
\usepackage{fancyhdr}

\pagestyle{fancy}
\fancyhead[L]{SteveROS Session Log}
\fancyhead[R]{2026-03-04}

\definecolor{codebg}{RGB}{245,245,245}
\definecolor{codeframe}{RGB}{200,200,200}
\lstset{
  backgroundcolor=\color{codebg},
  frame=single,
  rulecolor=\color{codeframe},
  basicstyle=\ttfamily\small,
  breaklines=true,
  columns=fullflexible,
  language=C++,
}

\title{SteveROS Development Session Log\\[0.5em]
\large MoveIt2 Integration, Hardware Driver Fixes, and Collision Configuration}
\author{SteveROS Team}
\date{March 4, 2026}

\begin{document}
\maketitle
\tableofcontents
\newpage

% ============================================================================
\section{Session Overview}

This document covers work done on the SteveROS humanoid robot platform on March 4, 2026.
The session focused on bringing up MoveIt2 motion planning with real Robstride motors
over CAN bus, diagnosing and fixing joint angle reporting issues, resolving collision
configuration problems, and investigating motor fault conditions during trajectory execution.

\subsection{System Configuration}
\begin{itemize}
  \item \textbf{OS:} Ubuntu 24.04, Linux 6.17
  \item \textbf{ROS Distribution:} ROS 2 Jazzy
  \item \textbf{Robot:} 20-DOF humanoid (kbot v2) with Robstride RS02/RS03/RS04 motors
  \item \textbf{Communication:} CAN bus (\texttt{can0}) at 1\,Mbps
  \item \textbf{Motors online:} 10/20 (both arms; legs not connected)
\end{itemize}

% ============================================================================
\section{Problem 1: MoveIt START\_STATE\_INVALID}

\subsection{Symptom}
When launching MoveIt2 with real hardware, motion planning immediately failed with:

\begin{lstlisting}[language={}]
CheckStartStateBounds: Joint 'dof_left_wrist_00' outside bounds
  by 6.02135 rad. Range: [-1.74533, 1.74533]
CheckStartStateBounds: Joint 'dof_right_shoulder_yaw_02' outside
  bounds by -1.68943 rad. Range: [-1.65806, 1.65806]
\end{lstlisting}

MoveIt rejected the robot's current joint state as invalid before any planning could begin.

\subsection{Root Cause Analysis}
The hardware driver's \texttt{motor\_to\_joint()} function converts raw motor encoder
readings to joint-space angles:

\begin{lstlisting}
// Original code (steveros_hardware.cpp)
double SteveROSHardware::motor_to_joint(
    const JointConfig & cfg, double motor_rad) const
{
  return cfg.sign * (motor_rad - cfg.zero_offset_rad);
}
\end{lstlisting}

The Robstride motor encoders report positions in the range $[-4\pi, 4\pi]$ (16-bit encoding).
When a motor is near the zero-offset boundary, the subtraction
\texttt{motor\_rad - zero\_offset\_rad} can produce values like $+6.02$\,rad or
$-1.69$\,rad---essentially $\pm 2\pi$ off from the true joint angle. These values
exceed the URDF joint limits, causing MoveIt to reject the start state.

\textbf{Key insight from testing:} The robot does not always start at the home position.
The arms can be in arbitrary poses when powered on, so the fix must handle any starting
configuration, not just recalibrate zero offsets.

\subsection{Fix Applied}

Added angle wrapping using \texttt{std::remainder()} to normalize the output to $[-\pi, \pi]$:

\begin{lstlisting}
// Fixed code (steveros_hardware.cpp)
double SteveROSHardware::motor_to_joint(
    const JointConfig & cfg, double motor_rad) const
{
  // Wrap to [-pi, pi] so joints near the zero-offset
  // boundary don't jump to +/-2pi.
  return std::remainder(
    cfg.sign * (motor_rad - cfg.zero_offset_rad),
    2.0 * M_PI);
}
\end{lstlisting}

\texttt{std::remainder(x, 2\pi)} returns the IEEE remainder, which is always in
$[-\pi, \pi]$. This ensures that regardless of how many full rotations the encoder
has accumulated, the reported joint angle stays within the physically meaningful range.

\subsection{File Modified}
\begin{itemize}
  \item \texttt{steveros\_hardware/src/steveros\_hardware.cpp} --- \texttt{motor\_to\_joint()} function
\end{itemize}

% ============================================================================
\section{Problem 2: Right Arm Collision False Positive}

\subsection{Symptom}
After fixing the start state bounds issue, the left arm could plan and execute
trajectories, but the right arm failed with:

\begin{lstlisting}[language={}]
CheckStartStateCollision failed: 1 contact(s) detected:
  KB_C_501X_Right_Bayonet_Adapter_Hard_Stop - RS03_4
\end{lstlisting}

MoveIt detected a collision between the right arm tip link
(\texttt{KB\_C\_501X\_Right\_Bayonet\_Adapter\_Hard\_Stop}) and the right hip motor link
(\texttt{RS03\_4}) at the current pose, preventing any planning for the right arm.

\subsection{Root Cause}
The SRDF (Semantic Robot Description Format) file defines which link pairs should
skip collision checking. The pair of the right arm end-effector and right hip motor
was missing from the \texttt{disable\_collisions} list. These links are on entirely
different kinematic chains (arm vs.\ leg) and are physically far apart in normal
operation, but their collision meshes can overlap at certain arm poses.

\subsection{Fix Applied}
Added \texttt{disable\_collisions} entries for arm tip and forearm links versus all
hip/leg links:

\begin{lstlisting}[language=XML]
<!-- Arm tips vs leg links -->
<disable_collisions
  link1="KB_C_501X_Right_Bayonet_Adapter_Hard_Stop"
  link2="RS03_4" reason="Never"/>
<disable_collisions
  link1="KB_C_501X_Right_Bayonet_Adapter_Hard_Stop"
  link2="RS03_5" reason="Never"/>
<!-- ... and similar for left arm, forearm links,
     and hip yoke links (16 entries total) -->
\end{lstlisting}

A total of 16 new collision pairs were disabled, covering both arms' tip and forearm
links against both legs' hip and upper leg links.

\subsection{File Modified}
\begin{itemize}
  \item \texttt{steveros\_moveit\_config/config/steveros.srdf} --- added 16 \texttt{disable\_collisions} entries
\end{itemize}

% ============================================================================
\section{Problem 3: Motor Faults During Trajectory Execution}

\subsection{Symptom}
After successfully planning and sending a trajectory to both arms, the shoulder motors
entered a fault state and stopped responding:

\begin{lstlisting}[language={}]
Motor 21 ('dof_right_shoulder_pitch_03') fault:
  fault_bits=0x01, mode=0, temp=38.0C
Motor 22 ('dof_right_shoulder_roll_03') fault:
  fault_bits=0x01, mode=0, temp=41.0C
Motor 11 ('dof_left_shoulder_pitch_03') fault:
  fault_bits=0x01, mode=0, temp=38.0C
Motor 12 ('dof_left_shoulder_roll_03') fault:
  fault_bits=0x01, mode=0, temp=41.0C
\end{lstlisting}

MoveIt reported trajectory execution timed out, and the motors remained in fault state.

\subsection{Analysis}
\begin{itemize}
  \item \texttt{fault\_bits=0x01}: Bit 0 of the Robstride fault register (overcurrent/stall protection)
  \item \texttt{mode=0}: Motor in \textbf{Reset} mode (not Run mode), meaning the motor
    firmware tripped the motor out of active control
  \item Temperatures were normal (38--41\textdegree C), ruling out thermal shutdown
  \item The faults began immediately after MoveIt sent the trajectory, suggesting the
    commanded motion was too aggressive for the shoulder motors (RS03, the largest actuators)
  \item The trajectory execution upper bound was $\sim$9.7\,s for the first attempt and
    $\sim$17.2\,s for the second, both of which timed out
\end{itemize}

\subsection{Current Status}
This issue remains under investigation. The hardware driver currently logs faults but
does not attempt automatic recovery (clear fault + re-enable). Potential solutions include:

\begin{enumerate}
  \item Reducing MoveIt trajectory velocity/acceleration scaling for shoulder joints
  \item Adding automatic fault recovery in the hardware driver's \texttt{read()} loop
  \item Tuning the motor PD gains (\texttt{kp}, \texttt{kd}) for the shoulder joints
  \item Investigating whether the trajectory waypoints cause sudden large torque demands
\end{enumerate}

% ============================================================================
\section{Additional Change: MoveIt Home State}

Added a \texttt{home} group state for the \texttt{both\_arms} planning group in the SRDF,
with all 10 arm joints set to 0. This makes ``home'' available in the MoveIt RViz
MotionPlanning plugin's Goal State dropdown for all three planning groups
(\texttt{right\_arm}, \texttt{left\_arm}, \texttt{both\_arms}).

\subsection{File Modified}
\begin{itemize}
  \item \texttt{steveros\_moveit\_config/config/steveros.srdf} --- added \texttt{both\_arms} home state
\end{itemize}

% ============================================================================
\section{Summary of All Changes}

\begin{table}[h]
\centering
\begin{tabular}{@{}lll@{}}
\toprule
\textbf{File} & \textbf{Change} & \textbf{Problem Solved} \\
\midrule
\texttt{steveros\_hardware/src/} & Added \texttt{std::remainder()} & START\_STATE\_INVALID \\
\texttt{steveros\_hardware.cpp} & wrapping in \texttt{motor\_to\_joint()} & (joint angles $\pm 2\pi$ off) \\
\midrule
\texttt{steveros\_moveit\_config/} & Added 16 collision disable & Right arm planning blocked \\
\texttt{config/steveros.srdf} & entries for arm vs.\ leg links & by false collision \\
\midrule
\texttt{steveros\_moveit\_config/} & Added \texttt{both\_arms} home & Home state not selectable \\
\texttt{config/steveros.srdf} & group state & for both\_arms group \\
\bottomrule
\end{tabular}
\end{table}

% ============================================================================
\section{Open Issues}

\begin{enumerate}
  \item \textbf{Shoulder motor faults:} Motors 11, 12, 21, 22 fault with \texttt{0x01}
    during trajectory execution. Needs trajectory tuning or fault recovery logic.
  \item \textbf{Leg motors offline:} Only 10/20 motors come online (arms only).
    Leg CAN connections not verified.
  \item \textbf{Trajectory timeout:} MoveIt's allowed execution duration may need
    adjustment for large motions, or trajectory velocity scaling should be reduced.
\end{enumerate}

\end{document}
